\section{Literature Review}
\label{sec:literature-review}

The literature separates into generation mechanisms, deployment
transformations, and comparative evaluation; the unqualified label
\emph{vocoder efficiency} conflates model structure, executable
representation, and measurement protocol.

\subsection{Generation Mechanisms}

WaveNet factorizes waveform likelihood sample by sample
\citep{oord2016wavenet}; LPCNet reduces learned synthesis through linear
prediction but retains a recurrent excitation generator
\citep{valin2019lpcnet}.  Adversarial vocoders instead use parallel
generators.  MelGAN combines multi-scale discrimination with feature matching
\citep{kumar2019melgan}; HiFi-GAN adds multi-period discrimination and
quality-footprint variants \citep{kong2020hifigan}; BigVGAN introduces
periodic activations and anti-aliased representations
\citep{lee2023bigvgan}.  These mechanisms differ in critical-path
parallelism, so capacity alone cannot predict latency.

iSTFT vocoders perform learned prediction near the acoustic-frame rate and
delegate waveform reconstruction to overlap-add; iSTFTNet established the
pattern by replacing a time-domain generator's final upsampling stages
with the inverse transform \citep{kaneko2022istftnet}.  APNet2 predicts amplitude
and phase spectra \citep{du2023apnet2}; Vocos predicts complex coefficients
with a constant-resolution ConvNeXt stack \citep{siuzdak2024vocos}; FreeV
uses a pseudo-inverse mel prior \citep{lv2024freev}; and RNDVoC separates
analytical range and learned null-space components \citep{li2025rndvoc}.
RFWave instead integrates a rectified flow in a multi-band spectral
representation \citep{liu2025rfwave}, making solver evaluations an explicit
latency-quality control.

VocosFormer is not a published architecture: the project inserts
residual-convolution and self-attention components into a capacity-matched
Vocos chassis, adopting them from the WavTokenizer decoder lineage
\citep{ji2025wavtokenizer}; the codec is cited for component provenance
only.

\subsection{Deployment Transformations}

Pruning changes both numerical weights and, only when supported, their
storage and execution representation.  Classical train-prune-retrain work
couples sparsity to compressed storage \citep{han2015pruning};
dense zero masks establish tolerance to weight removal but not
acceleration.
Quantization likewise specifies weights, activations, scales, zero-points,
covered operators, and backend \citep{jacob2018quantization}.  CPU dynamic
quantization, ONNX QDQ graphs, and GPU weight-only quantization are therefore
distinct treatments despite sharing an integer bit width.  Compilation
changes the captured executable graph; graph breaks, specialization mode,
operator support, and hardware determine whether TorchDynamo/TorchInductor
produces a gain \citep{ansel2024pytorch2}.  The vocoder literature itself
pursues efficiency through architecture and training: WaveRNN prunes its
recurrent weights during training to reach real-time mobile-CPU synthesis
\citep{kalchbrenner2018wavernn}, and WaveFM distills a flow-matching
vocoder toward single-step synthesis \citep{luo2025wavefm}; the
transformations above instead operate on fixed trained checkpoints.
Reproducible deployment
evidence must therefore state transformation scope, backend, operator
coverage, and the paired pre-transformation measurement, not transformation
names alone.

\subsection{Comparative Evaluation}

PESQ models perceived speech distortion \citep{rix2001pesq}, STOI estimates
intelligibility \citep{taal2011stoi}, and UTMOS predicts mean opinion scores
\citep{saeki2022utmos}; agreement among these proxies is still not a listening
test.  Shared-protocol vocoder comparisons improve internal validity while
remaining sensitive to implementation and conditioning choices
\citep{govalkar2019comparison}, and VocBench trains and evaluates six
vocoders in one shared environment with objective and subjective scoring
\citep{albadawy2022vocbench}.  Deployment-focused work shows that
batch speed does not establish low-latency streaming performance because
causality, chunk size, lookahead, loading overhead, and hardware can reverse
rankings \citep{yoneyama2025comparative}.  VOCODE does not claim streaming
evaluation or listener scoring.  Relative to VocBench it contributes twelve configurations trained from
scratch under recorded architecture-specific budgets, explicit checkpoint
provenance, inclusion and exclusion registers for executable deployment
artifacts, paired before-and-after measurements on one hardware profile,
and constraint-aware selection.
